\documentclass[10pt,a4paper]{article}
\usepackage[a4paper,margin=2cm]{geometry}
\usepackage{mathptmx}
\usepackage[T1]{fontenc}
\usepackage{booktabs}
\usepackage{enumitem}
\usepackage{titlesec}
\usepackage[hidelinks]{hyperref}
\urlstyle{same}
\usepackage{microtype}

\titleformat{\section}{\normalsize\bfseries}{\thesection.}{0.5em}{}
\titlespacing*{\section}{0pt}{7pt}{3pt}
\setlength{\parindent}{0pt}
\setlength{\parskip}{3pt}
\setlist{nosep,leftmargin=1.4em,topsep=1pt}
\pagestyle{empty}

\hypersetup{pdftitle={Time-Restricted Evidence Retrieval for Predicting Neuroscience Experimental Results with LLMs},
            pdfauthor={Md. Imtiaj Alam Sajin}}

\begin{document}

\begin{center}
{\large\bfseries Time-Restricted Evidence Retrieval for Predicting\\ Neuroscience Experimental Results with LLMs}\\[5pt]
Md. Imtiaj Alam Sajin (26-94090-2)\\[1pt]
CSC5515 Natural Language Processing, Initial Proposal\\[1pt]
Group G24 \quad Project type C (retrieval-augmented generation)\\[1pt]
American International University-Bangladesh, Fall 2026--27
\end{center}

\section{Problem and Motivation}
Scientists often predict the result of a study before running it, based on what earlier studies found. Luo et al.~[1] showed that large language models (LLMs) can do this well. In their benchmark, BrainBench, a model sees two versions of a \emph{Journal of Neuroscience} abstract that are identical except for the results, and it must pick the published one. LLMs were more accurate than human experts on this task. However, the models decide only from what they learned during training. Their choice is not linked to any evidence, and they cannot use studies they rarely saw. The authors suggested retrieval as a next step but did not test it. In this project I add retrieval of prior literature to this task. The retriever may only return papers published before the test article, so the model has to predict from earlier evidence and cannot simply find the answer.

\section{Research Questions}
\begin{itemize}
\item[RQ1.] Does adding time-restricted prior abstracts to the context improve accuracy compared with scoring without retrieval?
\item[RQ2.] How does the result change with the retriever (BM25 or dense) and with the number of retrieved abstracts?
\item[RQ3.] How much higher would accuracy be if the time restriction were removed? This measures the effect of leakage.
\end{itemize}
I expect retrieval to help most for topics that are well covered in earlier literature.

\section{Related Work}
Luo et al.~[1] introduced BrainBench and BrainGPT, a Llama-2 model tuned with LoRA on neuroscience literature from 2002--2022. Both decide by perplexity alone; this work adds dated evidence at test time. Luo et al.~[2] trained small GPT-2 models from scratch on neuroscience text and reached expert-level accuracy, also without retrieval. Lewis et al.~[3] introduced retrieval-augmented generation for question answering, where retrieval needs no time limit. Halawi et al.~[4] forecast real-world events using news retrieved before a cutoff date; here the target is scientific results, and the decision is made by comparing likelihoods rather than by generating a forecast. BackTrend~[5] uses date-filtered retrieval to predict research trends in AI and machine learning, which is a different domain and task.

\section{Data}
\textbf{Test set.} The 200 human-written BrainBench items, built from 2023 \emph{Journal of Neuroscience} abstracts. I inspected them in the authors' experiment repository~[6]. They cover five subfields: Behavioral/Cognitive (75), Systems/Circuits (51), Neurobiology of Disease (25), Development/Plasticity/Repair (24) and Cellular/Molecular (25). Each item changes between one and six result sentences (45, 77, 41, 24, 8 and 5 items respectively), and the result part has 65 words on average. The label is binary (original or altered); I randomise the order of the two versions so the classes are balanced.

\textbf{Development set.} The 100 BrainBench items generated by GPT-4 and checked by experts~[1]. They are used only to choose settings such as the number of retrieved abstracts. The test set is used only for final results.

\textbf{Retrieval corpus.} 3,451 \emph{Journal of Neuroscience} abstracts from 2017--2022, released with the BrainBench code~[6] (320 words on average), all published before the test articles. If more coverage is needed, PubMed abstracts published before 2023 will be added through the NCBI E-utilities.

\textbf{Licence.} The BrainBench code is released under Apache 2.0. The test-case files are distributed on Hugging Face by the BrainGPT group, and their licence will be reported in the final proposal. PubMed abstracts are used only for retrieval and are not redistributed.

Table~\ref{tab:samples} shows five sample items. Only the changed words are listed.

\begin{table}[h]
\centering\small
\caption{Five sample test items (original $\rightarrow$ altered).}
\label{tab:samples}
\begin{tabular}{lllp{6.2cm}}
\toprule
Item & Subfield & Result sentences & Change in the altered version\\
\midrule
BC-0   & Behavioral/Cognitive          & 4 & anterior $\rightarrow$ posterior (location of cortical patches)\\
SC-10  & Systems/Circuits              & 3 & reduces $\rightarrow$ increases\\
ND-39  & Neurobiology of Disease       & 3 & short-term $\rightarrow$ long-term\\
DP-24  & Development/Plasticity/Repair & 2 & phasic $\rightarrow$ tonic, and tonic $\rightarrow$ phasic\\
CM-28  & Cellular/Molecular            & 3 & knocking down $\rightarrow$ overexpressing\\
\bottomrule
\end{tabular}
\end{table}

\section{Models and Baselines}
All models run locally with Hugging Face Transformers on one 16\,GB GPU, with 4-bit quantisation where needed. The main model is Llama-2-7B~[7], whose pretraining data end in 2022, before the test abstracts were published. Mistral-7B-v0.1 is a second model, and the public BrainGPT LoRA adapter~[1] is a domain-tuned comparison.
\begin{itemize}
\item \emph{Non-LLM baseline:} BM25 support. The version whose result sentences better match the retrieved abstracts is chosen.
\item \emph{Zero-shot LLM baseline:} the original BrainBench rule. The version with lower perplexity is chosen, without retrieval. A prompted zero-shot variant that asks which version is real is also reported.
\item \emph{Main method:} a query is built from the background and method sentences, and the top-$k$ abstracts published before the test article are retrieved. They are placed before each version, and the version with lower perplexity is chosen. Confidence is the difference in log-perplexity.
\end{itemize}

\section{Evaluation}
The main metric is accuracy, because the task is a balanced two-way choice and accuracy is also used for experts and models in~[1]. I report 95\% bootstrap confidence intervals and use McNemar tests for paired comparisons. Calibration is measured as accuracy per confidence quintile and as the AUROC of confidence. Perplexity scoring is deterministic, so variance comes from bootstrap resampling and three query variants; the prompted baseline is run with three seeds, and unparseable answers are counted as wrong. The ablations vary the number of retrieved abstracts ($k \in \{0,1,3,5\}$), BM25 against dense retrieval, the time filter (on or off) and the model. For the error analysis, at least 50 errors of the main method are grouped by the type of change (direction, brain region, cell or molecule type, time scale) and by what the retriever returned (nothing relevant, supporting evidence, conflicting evidence).

\section{Risks and Backup Plan}
\begin{itemize}
\item If accuracy is above 75\% on the 200 test items, there will be fewer than 50 errors. In that case, errors from the development set and from the second model are added and reported separately.
\item Newer models may have seen the 2023 abstracts during training. The main model predates them, and the time-filter ablation measures leakage directly.
\item Retrieved abstracts are truncated so that the input fits the 4,096-token context of the main model.
\item If retrieval does not improve accuracy, a reranking step is added before scoring. RQ3 and the calibration analysis still give results either way.
\end{itemize}

\section{Task Split and Timeline}
This is an individual project, so all tasks are done by me.

\begin{center}\small
\begin{tabular}{lp{13.2cm}}
\toprule
Week & Task\\
\midrule
6      & Final proposal; data download, cleaning script and statistics; baselines on the development set\\
7--8   & Retrieval index (BM25 and dense); time-restricted retrieval and conditional scoring\\
9--10  & Test runs, ablations, confidence intervals and calibration\\
11     & Error analysis; code package with README, requirements and one script for the main table\\
12     & Final report in the course \LaTeX{} template\\
\bottomrule
\end{tabular}
\end{center}

\section*{References}
\small
\begin{enumerate}[label={[\arabic*]},leftmargin=2em]
\item X. Luo et al. Large language models surpass human experts in predicting neuroscience results. \emph{Nature Human Behaviour}, 9(2):305--315, 2025.
\item X. Luo et al. Matching domain experts by training from scratch on domain knowledge. arXiv:2405.09395, 2024.
\item P. Lewis et al. Retrieval-augmented generation for knowledge-intensive NLP tasks. In \emph{NeurIPS}, 2020.
\item D. Halawi et al. Approaching human-level forecasting with language models. In \emph{NeurIPS}, 2024.
\item BackTrend: Evaluating scientific weak-signal prediction via backward reconstruction. arXiv:2609.24921, 2026.
\item BrainBench experiment code. \url{https://github.com/braingpt-lovelab/brainbench_experiments}
\item H. Touvron et al. Llama 2: Open foundation and fine-tuned chat models. arXiv:2307.09288, 2023.
\end{enumerate}

\end{document}
